%% ----------------------------------------------------------------------------
%%  §7.6  Funzioni di tre o più variabili reali
%% ----------------------------------------------------------------------------
\section{Funzioni di tre o più variabili reali}
\label{sec:n-variabili}

Indichiamo con $\R^{3}$ l'insieme delle \emph{terne ordinate} di numeri reali e, più in
generale, con $\R^{n}$ l'insieme delle \emph{$n$-ple ordinate}:
\begin{equation}\label{eq:Rn}
\begin{gathered}
  \R^{3}=\set{(x,y,z) : x,y,z\in\R},\\
  \R^{n}=\set{(x_{1},x_{2},\dots,x_{n}) : x_{i}\in\R\ \perogni i=1,\dots,n}.
\end{gathered}
\end{equation}
Una legge $f$ che a ogni elemento di un insieme $D\subseteq\R^{n}$ fa corrispondere uno e
un solo numero reale è detta una \emph{funzione di $n$ variabili} e si denota con
$f(x_{1},x_{2},\dots,x_{n})$, oppure, nel caso $n=3$, con $f(x,y,z)$; l'insieme $D$ è il
\emph{dominio} di $f$. La distanza tra due punti di $\R^{n}$ è
$\sqrt{(x_{1}-y_{1})^{2}+\dots+(x_{n}-y_{n})^{2}}$, e con essa si definiscono gli intorni
(sferici, se $n=3$), i limiti e la continuità esattamente come nel piano. Vediamo più in
dettaglio le derivate parziali.

\begin{definizione}[Derivate parziali in $\R^{n}$]\label{def:parziali-n}
La \emph{derivata parziale} di $f$ rispetto alla variabile $x_{i}$ nel punto
$(x_{1},\dots,x_{n})$ è il limite
\begin{equation}\label{eq:parziale-n}
  \lim_{h\to0}\frac{f(x_{1},\dots,x_{i}+h,\dots,x_{n})-f(x_{1},\dots,x_{i},\dots,x_{n})}{h},
\end{equation}
purché esista e sia finito; si denota con uno dei simboli $f_{x_{i}}$,
$\frac{\partial f}{\partial x_{i}}$, $D_{x_{i}}f$, $D_{i}f$. Il vettore
$\nabla f=(f_{x_{1}},\dots,f_{x_{n}})$ è il \emph{gradiente} di $f$.
\end{definizione}

Anche qui si deriva rispetto a una variabile considerando costanti tutte le altre. Ad
esempio, le derivate parziali della funzione di tre variabili
\begin{equation}\label{eq:f-tre-variabili}
  f(x,y,z)=x^{2}+y^{2}+z^{2}+xyz
\end{equation}
sono $f_{x}=2x+yz$, $f_{y}=2y+xz$, $f_{z}=2z+xy$. La funzione di $n$ variabili detta
\emph{modulo} o \emph{norma},
\begin{equation}\label{eq:norma}
  f(x_{1},\dots,x_{n})=\sqrt{x_{1}^{2}+x_{2}^{2}+\dots+x_{n}^{2}},
\end{equation}
che misura la distanza del punto dall'origine, ammette in ogni punto diverso
dall'origine le derivate parziali
\begin{equation}\label{eq:norma-derivate}
  f_{x_{i}}=\frac{2x_{i}}{2\sqrt{x_{1}^{2}+\dots+x_{n}^{2}}}=\frac{x_{i}}{f},
  \qquad i=1,2,\dots,n,
\end{equation}
come si vede derivando la radice con la regola della catena. Il gradiente della norma è
quindi il versore $\frac{(x_{1},\dots,x_{n})}{f}$, diretto radialmente verso l'esterno:
la distanza dall'origine cresce più rapidamente allontanandosi in direzione radiale.

\begin{esempio}[La portanza come funzione di tre variabili]\label{es:portanza-tre}
La portanza $L(\rho,V,C_{L})=\frac12\rho SV^{2}C_{L}$ è una funzione di tre variabili, con
\[
  \frac{\partial L}{\partial\rho}=\frac12SV^{2}C_{L}=\frac L\rho, \qquad
  \frac{\partial L}{\partial V}=\rho SVC_{L}=\frac{2L}V, \qquad
  \frac{\partial L}{\partial C_{L}}=\frac{L}{C_{L}} .
\]
Moltiplicando ciascuna derivata per la variabile corrispondente e dividendo per $L$ si
ottengono le \emph{sensibilità relative} $1$, $2$, $1$: una variazione percentuale piccola
della velocità produce una variazione percentuale \emph{doppia} della portanza, mentre
le variazioni di densità e di $C_{L}$ si trasmettono in egual misura. Ad esempio, salendo
di quota a velocità e assetto costanti, la portanza diminuisce nella stessa proporzione
della densità; per ristabilire l'equilibrio occorre aumentare $C_{L}$ oppure la velocità,
e per quest'ultima basta un aumento percentuale pari a circa la metà. Daremo un
fondamento preciso a questi ragionamenti con il differenziale, nel
\cref{sec:differenziabilita}.
\end{esempio}

Come per le funzioni di due variabili, si definiscono le \emph{derivate parziali
seconde} $f_{x_{i}x_{j}}=(f_{x_{i}})_{x_{j}}$, che vengono a costituire la seguente
matrice quadrata $n\times n$, detta \emph{matrice hessiana}:
\begin{equation}\label{eq:matrice-hessiana}
  \begin{pmatrix}
    f_{x_{1}x_{1}} & f_{x_{1}x_{2}} & \cdots & f_{x_{1}x_{n}}\\
    f_{x_{2}x_{1}} & f_{x_{2}x_{2}} & \cdots & f_{x_{2}x_{n}}\\
    \vdots & \vdots & \ddots & \vdots\\
    f_{x_{n}x_{1}} & f_{x_{n}x_{2}} & \cdots & f_{x_{n}x_{n}}
  \end{pmatrix}.
\end{equation}
Applicando il teorema di Schwarz alla funzione delle sole due variabili $x_{i}$, $x_{j}$
che si ottiene lasciando fissate tutte le altre, con il ruolo di parametri, si ottiene
\begin{equation}\label{eq:schwarz-n}
  f_{x_{i}x_{j}}=f_{x_{j}x_{i}} \qquad\perogni i,j=1,\dots,n,
\end{equation}
purché tali derivate seconde esistano e siano continue. Pertanto la matrice hessiana è
\emph{simmetrica} nel caso in cui tutte le derivate parziali seconde siano continue. Con
le stesse definizioni e la stessa dimostrazione del \cref{teo:fermat-due} si prova che in
un punto di massimo o di minimo relativo interno al dominio risulta
\begin{equation}\label{eq:fermat-n}
  f_{x_{1}}=f_{x_{2}}=\dots=f_{x_{n}}=0,
\end{equation}
nell'ipotesi che le derivate parziali prime esistano nel punto.

\begin{osservazione}[La condizione sufficiente in $n$ variabili]\label{oss:hessiana-n}
Una condizione sufficiente per i massimi e i minimi, in termini della matrice hessiana,
esiste ma è più complessa della formulazione data nel caso di due variabili. Per una
funzione con derivate seconde continue in un intorno del punto critico, essa si
esprime dicendo che la matrice hessiana nel punto critico è \emph{definita positiva}
(minimo) o \emph{definita negativa} (massimo), nozioni che si studiano nei corsi di
Geometria; se la forma quadratica associata assume valori di segno opposto lungo due
direzioni, il punto critico è di sella. Un criterio pratico, dovuto a Sylvester, afferma
che una matrice simmetrica è definita positiva se e solo se sono positivi i minori
principali «di nord-ovest», cioè i determinanti delle sottomatrici formate dalle prime
$k$ righe e colonne, per $k=1,\dots,n$; è definita negativa se e solo se tali minori
hanno segni alterni, a cominciare da un negativo. Per $n=2$ si ritrovano le condizioni
$f_{xx}>0$ (oppure $f_{xx}<0$) e $H>0$ del \cref{teo:hessiano}. Come vedremo nel
\cref{sec:differenziabilita} (per $n=2$; la dimostrazione è identica per $n$ qualsiasi),
la derivata seconda nel punto critico $P$ della restrizione di $f$ alla retta per $P$ di
direzione $v$ è la forma quadratica $v^{T}H(P)\,v$: è questo il motivo per cui, se essa
assume valori di segno opposto per due direzioni, $P$ è un punto di sella.
\end{osservazione}

Il programma seguente trova i punti critici della funzione \eqref{eq:f-tre-variabili} e,
per ciascuno, calcola la forma quadratica $v^{T}Hv$ per i $26$ vettori $v$ con componenti
in $\set{-1,0,1}$ non tutte nulle.

\begin{python}[Punti critici di una funzione di tre variabili]
import itertools
import sympy as sp

x, y, z = sp.symbols("x y z", real=True)
f = x**2 + y**2 + z**2 + x*y*z
grad = [sp.diff(f, v) for v in (x, y, z)]
H = sp.hessian(f, (x, y, z))
print("grad f =", grad)
# direzioni di prova v; la derivata seconda lungo v è v^T H v
dirs = [sp.Matrix(v) for v in itertools.product((-1, 0, 1), repeat=3)
        if any(v)]
for P in sp.solve(grad, [x, y, z], dict=True):
    q = [int((v.T * H.subs(P) * v)[0]) for v in dirs]
    c = ", ".join(f"{int(P[s]):2d}" for s in (x, y, z))
    print(f"P = ({c}):  min v^T H v = {min(q):2d},"
          f"  max v^T H v = {max(q)}")
\end{python}
\begin{risultato}
grad f = [2*x + y*z, x*z + 2*y, x*y + 2*z]
P = (-2, -2, -2):  min v^T H v = -6,  max v^T H v = 10
P = (-2,  2,  2):  min v^T H v = -6,  max v^T H v = 10
P = ( 0,  0,  0):  min v^T H v =  2,  max v^T H v = 6
P = ( 2, -2,  2):  min v^T H v = -6,  max v^T H v = 10
P = ( 2,  2, -2):  min v^T H v = -6,  max v^T H v = 10
\end{risultato}

Nell'origine la matrice hessiana è $2I$, il doppio della matrice identità, e
$v^{T}Hv=2(v_{1}^{2}+v_{2}^{2}+v_{3}^{2})>0$ per ogni $v\neq0$: è definita positiva, e
l'origine è un punto di minimo relativo (come si vede anche osservando che, vicino
all'origine, il termine cubico $xyz$ è trascurabile rispetto a $x^{2}+y^{2}+z^{2}$). Negli
altri quattro punti critici, del tipo $(\pm2,\pm2,\pm2)$ con prodotto delle coordinate
negativo, la forma quadratica assume valori di segno opposto: sono punti di sella.
